\section{Finite ring maps}
\label{section-finite}

\noindent
Here is the definition.

\begin{definition}
\label{definition-finite-ring-map}
Let $\varphi : R \to S$ be a ring map. We say $\varphi : R \to S$ is
{\it finite} if $S$ is finite as an $R$-module.
\end{definition}

\begin{lemma}
\label{lemma-finite-module-over-finite-extension}
Let $R \to S$ be a finite ring map.
Let $M$ be an $S$-module.
Then $M$ is finite as an $R$-module if and only if $M$ is finite
as an $S$-module.
\end{lemma}

\begin{proof}
One of the implications follows from
Lemma \ref{lemma-finite-over-subring}.
To see the other assume that $M$ is finite as an $S$-module.
Pick $x_1, \ldots, x_n \in S$ which generate $S$ as an $R$-module.
Pick $y_1, \ldots, y_m \in M$ which generate $M$ as an $S$-module.
Then $x_i y_j$ generate $M$ as an $R$-module.
\end{proof}

\begin{lemma}
\label{lemma-finite-transitive}
Suppose that $R \to S$ and $S \to T$ are finite ring maps.
Then $R \to T$ is finite.
\end{lemma}

\begin{proof}
If $t_i$ generate $T$ as an $S$-module and $s_j$ generate $S$ as an
$R$-module, then $t_i s_j$ generate $T$ as an $R$-module.
(Also follows from
Lemma \ref{lemma-finite-module-over-finite-extension}.)
\end{proof}

\begin{lemma}
\label{lemma-finite-finite-type}
Let $\varphi : R \to S$ be a ring map.
\begin{enumerate}
\item If $\varphi$ is finite, then $\varphi$ is of finite type.
\item If $S$ is of finite presentation as an $R$-module, then
$\varphi$ is of finite presentation.
\end{enumerate}
\end{lemma}

\begin{proof}
For (1) if $x_1, \ldots, x_n \in S$ generate $S$ as an $R$-module,
then $x_1, \ldots, x_n$ generate $S$ as an $R$-algebra. For (2),
suppose that $\sum r_j^ix_i = 0$, $j = 1, \ldots, m$ is a set
of generators of the relations among the $x_i$ when viewed as
$R$-module generators of $S$. Furthermore, write
$1 = \sum r_ix_i$ for some $r_i \in R$ and
$x_ix_j = \sum r_{ij}^k x_k$ for some $r_{ij}^k \in R$.
Then
$$
S = R[t_1, \ldots, t_n]/
(\sum r_j^it_i,\ 1 - \sum r_it_i,\ t_it_j - \sum r_{ij}^k t_k)
$$
as an $R$-algebra which proves (2).
\end{proof}

\noindent
For more information on finite ring maps, please see
Section \ref{section-finite-ring-extensions}.


